% Figure: the relative-motion transfer function Z/Y of a base-excited
% seismic transducer, showing the seismometer region (r >> 1, Z/Y -> 1)
% and the accelerometer region (r << 1, Z/Y -> r^2, flat in
% acceleration). Z/Y = r^2 / sqrt((1-r^2)^2 + (2 zeta r)^2).
\begin{figure}[htbp]
\centering
\begin{tikzpicture}
\begin{axis}[
  width=12cm, height=7.2cm,
  xmode=log, ymode=log,
  xlabel={frequency ratio $r=\omega/\omega_n$},
  ylabel={relative response $Z/Y$},
  xmin=0.05, xmax=20, ymin=0.001, ymax=20,
  axis lines=left,
  domain=0.05:20, samples=300,
  legend pos=north west,
  legend cell align=left,
  every axis x label/.style={at={(current axis.right of origin)}, anchor=west},
  every axis y label/.style={at={(current axis.above origin)}, anchor=south},
  grid=major, grid style={enggray!20},
  tick label style={font=\scriptsize},
  label style={font=\small},
  legend style={font=\scriptsize},
]
% zeta = 0.7, near-flat accelerometer region.
\addplot[engblue, very thick] {x^2/sqrt((1-x^2)^2 + (2*0.7*x)^2)};
\addlegendentry{$\zeta=0.7$}
% zeta = 0.05, lightly damped (peaks at resonance).
\addplot[engred, very thick, dashed] {x^2/sqrt((1-x^2)^2 + (2*0.05*x)^2)};
\addlegendentry{$\zeta=0.05$}
% low-frequency accelerometer asymptote Z/Y = r^2
\addplot[enggray, thin, forget plot, domain=0.05:0.6] {x^2};
% high-frequency seismometer asymptote Z/Y = 1
\addplot[enggray, thin, forget plot, domain=2:20] {1};
\node[engblack, font=\scriptsize, anchor=south] at (axis cs:0.13,0.04) {accelerometer ($Z/Y\to r^2$)};
\node[engblack, font=\scriptsize, anchor=south] at (axis cs:6,1.2) {seismometer ($Z/Y\to 1$)};
\end{axis}
\end{tikzpicture}
\caption[Seismic-transducer design regions]{Relative displacement
$Z=x-y$ of a seismic mass across its spring, per unit base motion $Y$,
for a base excited at frequency ratio $r=\omega/\omega_n$. Two usable
regions sit at the ends. Well below resonance ($r\ll1$) the relative
motion follows $Z/Y\to r^2=\omega^2/\omega_n^2$, proportional to base
acceleration, so a stiff high-frequency instrument reads acceleration:
the accelerometer region. Well above resonance ($r\gg1$) the relative
motion follows $Z/Y\to1$, equal to the base displacement, because the
seismic mass stands still in space while the case moves around it: the
seismometer region. The lightly damped curve (red dashed, $\zeta=0.05$)
peaks sharply at resonance and has only narrow flat regions; the
$\zeta=0.7$ curve (blue) maximises the flat accelerometer band, the
reason practical accelerometers are damped near $0.7$ of critical.
Computed from $Z/Y=r^2/[(1-r^2)^2+(2\zeta r)^2]^{1/2}$.}
\label{fig:vol03ch06:seismometer}
\end{figure}
